\unit{Exposé 232 -- Picard schemes: existence theorems}
\sid{232}
\src{143}
\seminar{14th year, 1961/62, no. 232\\February 1962}
\paperhead{DESCENT TECHNIQUE AND EXISTENCE THEOREMS IN ALGEBRAIC GEOMETRY\\
V. PICARD SCHEMES: EXISTENCE THEOREMS}%
  {by Alexander GROTHENDIECK}

\fsection{1}{Relative Picard groups and functors}

For every prescheme (more generally, every ringed space) $X$, we call the group
of isomorphism classes of invertible Modules (that is, Modules locally
isomorphic to $\mathcal O_X$) on $X$ the (absolute) \emph{Picard group} of $X$,
and denote it by $\operatorname{Pic}(X)$. Thus there is a canonical isomorphism
\begin{equation}
\operatorname{Pic}(X)\xrightarrow{\sim}H^1(X,\mathcal O_X^*),
\tag{1.1}
\end{equation}
where $\mathcal O_X^*$ denotes the sheaf of units of $\mathcal O_X$ (which is
indeed identified with the sheaf of automorphisms of the standard invertible
Module $\mathcal O_X$). Clearly $X\longmapsto\operatorname{Pic}(X)$ is a
contravariant functor in $X$, and isomorphism (1.1) is functorial.

If $X$ is a prescheme over another prescheme $S$, then as $S'$ varies in the
category $(\mathrm{Sch})/S$ of preschemes over $S$, the preceding construction
gives a contravariant functor
$S'\longmapsto\operatorname{Pic}(X\times_SS')$. This functor has no chance of
being “representable” ([4], II, A), because the nontrivial automorphisms of the
invertible Modules to be classified prevent it from being “local in nature”
([5], IV, 5.4). It must therefore be “made local.” In general, for every
relative prescheme $X/S$, introduce the relative group
\begin{equation}
\operatorname{Pic}'(X/S)=H^0\bigl(S,R^1f_*(\mathcal O_X^*)\bigr),
\tag{1.2}
\end{equation}
where $f:X\longrightarrow S$ is the structural morphism (compare [4], II,
C 3). In the cited passage this group is called the relative Picard group; for
reasons that will appear, it is preferable here to call it the
\emph{restricted relative Picard group} of $X/S$. As $S'$ varies in
$(\mathrm{Sch})/S$, the assignment
$S'\longmapsto\operatorname{Pic}'(X\times_SS'/S')$ is a contravariant functor
in $S'$, also denoted $\operatorname{Pic}'_{X/S}$ and therefore defined
essentially by
\begin{equation}
\operatorname{Pic}'_{X/S}(S')=
\operatorname{Pic}'(X\times_SS'/S').
\tag{1.3}
\end{equation}

\src{144}
This functor is now “local in nature,” since the needed adjustment has been
made. Intuitively, the right-hand side of (1.3) is the set of “algebraic
families” of classes of invertible sheaves on (the fibers of) $X/S$, indexed
by the parameter prescheme $S'/S$. When the functor $\operatorname{Pic}'$ is
representable, the prescheme over $S$ representing it is denoted
$\operatorname{Pic}'_{X/S}$ and called the \emph{Picard prescheme} of $X$ over
$S$; thus one then has
\begin{equation}
\operatorname{Hom}_S(S',\operatorname{Pic}'_{X/S})
\simeq\operatorname{Pic}'_{X/S}(S')
=\operatorname{Pic}'(X\times_SS'/S').
\tag{1.4}
\end{equation}

There are, however, important cases in which $\operatorname{Pic}'_{X/S}$ is not
representable (for example, a “Brauer--Severi” variety over a field $k$ with
no $k$-rational point), although a relative Picard prescheme still has a
natural definition. The reason is that the definition of
$\operatorname{Pic}'$ from the absolute Picard groups
$\operatorname{Pic}(X\times_SS'/S')$ has not yet been localized enough: more
precisely, $\operatorname{Pic}'$ is not in general “compatible with faithfully
flat descent.” Let us make this explicit.

Let $\mathfrak M$ be the class of prescheme morphisms that are
\emph{faithfully flat and quasi-compact}; this class is stable under base
change and composition. Let $P$ be a contravariant functor from
$(\mathrm{Sch})/S$ to the category of sets, and for every $S$-morphism
$u:T'\longrightarrow T$ with $u\in\mathfrak M$, consider the diagram
\begin{equation}
P(T)\longrightarrow P(T')\rightrightarrows P(T'\times_TT')
\tag{1.5}
\end{equation}
obtained by applying $P$ to
\[
T\longleftarrow T'\mathrel{\substack{\xleftarrow{\,\mathrm{pr}_1\,}\\[-0.35ex]
\xleftarrow[\,\mathrm{pr}_2\,]{}}}T'\times_TT'.
\]
If $P$ is representable, descent theory ([4], I, B, Th. 2) implies that (1.5)
is exact for every $u\in\mathfrak M$. This is expressed by saying that $P$ is
compatible with $\mathfrak M$, in this case that $P$ is “compatible with
faithfully flat descent,” or that the “presheaf” $P$ on $(\mathrm{Sch})/S$ is
a “sheaf” for the notion of localization supplied by $\mathfrak M$. For an
arbitrary $P$, a standard procedure, familiar from ordinary topological
localization, associates with it a “sheaf” $\mathcal P$ and a functor
homomorphism $P\longrightarrow\mathcal P$ universal in the evident sense. The
calculation of $\mathcal P$ can be made explicit as follows. To define
$\mathcal P(T)$, for every $T'$ over $T$ such that
$u:T'\longrightarrow T$ belongs to $\mathfrak M$, let

\src{145}
$\overline H^0(T'/T,P)$ be the subset of $P(T')$ consisting of elements $\xi$
whose two images $\xi_1,\xi_2$ in $P(T'\times_TT')$ have the same image in
$P(T'')$ for some morphism $v:T''\longrightarrow T'\times_TT'$ with
$v\in\mathfrak M$. [N.B. -- The set $\overline H^0$ so defined is therefore
larger than the set $H^0(T'/T,P)$ introduced in [4], I, A, 4(a).] As $T'$ varies
over a fixed $T$ (still with $u\in\mathfrak M$), the
$\overline H^0(T'/T,P)$ form an inductive system when the set of $T'$ is given
the preorder defined by domination, and one sets
\begin{equation}
\mathcal P(T)=\varinjlim_{T'}\overline H^0(T'/T,P).
\tag{1.6}
\end{equation}
The functorial law in $T$ for this expression is immediate.

When
\[
P(T)=\operatorname{Pic}(X\times_ST),
\]
the contravariant functor on $(\mathrm{Sch})/S$ defined by (1.6) is called the
\emph{relative Picard functor} of $X$ over $S$ and denoted
$\operatorname{Pic}_{X/S}$; the group $\operatorname{Pic}_{X/S}(S)$ is called
the \emph{relative Picard group} of $X$ over $S$ and denoted
$\operatorname{Pic}(X/S)$. There is then an evident bijection
\begin{equation}
\operatorname{Pic}_{X/S}(T)\simeq\operatorname{Pic}(X\times_ST/T).
\tag{1.7}
\end{equation}

Thus an element of $\operatorname{Pic}(X/S)$ is defined by an element $\xi'$ of
a group $\operatorname{Pic}(X\times_SS')$, where $S'\longrightarrow S$ is
faithfully flat and quasi-compact, such that there is a faithfully flat and
quasi-compact morphism $S''\longrightarrow S'\times_SS'$ for which the two
inverse images of $\xi'$ in $\operatorname{Pic}(X\times_SS'')$ coincide. An
element $\xi'$ of $\operatorname{Pic}(X\times_SS')$ and an element $\xi_1$ of
$\operatorname{Pic}(X\times_SS_1)$ satisfying the conditions just stated
define the same element of $\operatorname{Pic}(X/S)$ if and only if there is a
faithfully flat and quasi-compact morphism
$S'_1\longrightarrow S'\times_SS_1$ such that the images of the two elements
in $\operatorname{Pic}(X\times_SS'_1)$ are equal. It is often convenient still
to work with the functor $P'=\operatorname{Pic}'_{X/S}$ introduced above; one
checks at once that the canonical morphism $P\longrightarrow P'$ induces an
isomorphism
\begin{equation}
\mathcal P\xrightarrow{\sim}\mathcal P',
\tag{1.8}
\end{equation}
which gives an often more convenient description of
$\operatorname{Pic}_{X/S}$ in terms of $\operatorname{Pic}'_{X/S}=P'$. By 2.3
below, when $P$ is replaced by $P'$ in the description of
$\operatorname{Pic}(X/S)$ just given, one may in fact
\src{146}
take $S''=S'\times_SS'$ and $S'_1=S'\times_SS_1$, at least under the conditions
specified in the cited passage.

When the functor $\operatorname{Pic}_{X/S}$ is representable, one says that
$X/S$ admits a Picard prescheme, and the prescheme over $S$ representing the
functor is called the \emph{Picard prescheme} of $X$ over $S$, still denoted
$\operatorname{Pic}_{X/S}$. It plainly suffices that
$P'=\operatorname{Pic}'_{X/S}$ be representable, for then $P'$ is already a
“sheaf,” and (1.8) shows that $P'\longrightarrow\mathcal P'$ is the canonical
morphism
\begin{equation}
\operatorname{Pic}'_{X/S}\longrightarrow\operatorname{Pic}_{X/S},
\tag{1.9}
\end{equation}
which is then an isomorphism. Thus our terminology agrees with that introduced
above in (1.4). In general, when $\operatorname{Pic}_{X/S}$ exists, it is
defined by the functorial isomorphism
\begin{equation}
\operatorname{Hom}_S(S',\operatorname{Pic}_{X/S})\xrightarrow{\sim}
\operatorname{Pic}(X\times_SS'/S').
\tag{1.10}
\end{equation}

\fsection{2}{Relations among the various relative and absolute Picard groups}

\result{PROPOSITION}{2.1} Let $f:X\longrightarrow S$ be a morphism such that
$\mathcal O_S\xrightarrow{\sim}f_*(\mathcal O_X)$. Then there is an exact
sequence
\[
0\longrightarrow\operatorname{Pic}(S)\longrightarrow\operatorname{Pic}(X)
\longrightarrow\operatorname{Pic}'(X/S).
\]
When $X$ admits a section over $S$, the last morphism is surjective; that is,
there is an isomorphism
\[
\operatorname{Pic}'(X/S)\xrightarrow{\sim}
\operatorname{Pic}(X)/\operatorname{Pic}(S).
\]

The exact sequence may be regarded as the low-degree exact sequence associated
with the Leray spectral sequence for $f$ and $\mathcal O_X^*$. The second
assertion is likewise formal.

\result{PROPOSITION}{2.2} Let $f:X\longrightarrow S$ be a quasi-compact
separated morphism such that
$\mathcal O_S\xrightarrow{\sim}f_*(\mathcal O_X)$, and let
$S'\longrightarrow S$ be faithfully flat and quasi-compact. Then:

\noindent (i) $\operatorname{Pic}'(X/S)\longrightarrow
\operatorname{Pic}'(X\times_SS'/S')$ is injective;

\noindent (ii) if $X$ admits a section locally over $S$ (that is, every
$s\in S$ has an open neighborhood $U$ such that $X|U$ has a section over $U$),
then the diagram
\[
\operatorname{Pic}'(X/S)\longrightarrow
\operatorname{Pic}'(X\times_SS'/S')\rightrightarrows
\operatorname{Pic}'(X\times_SS''/S'')
\]

\src{147}
where $S''=S'\times_SS'$, is exact.

The first assertion follows from the elementary properties of faithfully flat
descent and the following general remark. If $f:X\longrightarrow S$ satisfies
$\mathcal O_S\xrightarrow{\sim}f_*(\mathcal O_X)$, then the functor
$\mathcal F\longmapsto f^*(\mathcal F)$ from locally free finite-type Modules
on $S$ to locally free finite-type Modules on $X$ is fully faithful, and its
essential image consists of the Modules $\mathcal G$ on $X$ such that
$f_*(\mathcal G)$ is locally free and the canonical homomorphism
\[
f^*\bigl(f_*(\mathcal G)\bigr)\longrightarrow\mathcal G
\]
is an isomorphism. The second assertion was proved by descent theory in [4],
I, B, 4.

The results of 2.2 may also be stated as follows.

\result{COROLLARY}{2.3} Under the conditions of 2.2, the canonical homomorphism
(1.9) $\operatorname{Pic}'_{X/S}\longrightarrow\operatorname{Pic}_{X/S}$ is
injective, and is bijective if $X$ admits a section locally over $S$. (Thus, in
the latter case, the relative Picard group $\operatorname{Pic}(X/S)$ and the
restricted relative Picard group $\operatorname{Pic}'(X/S)$ are identical.)

Combining this with 2.1 gives:

\result{COROLLARY}{2.4} Under the conditions of 2.2, there is an exact sequence
\[
0\longrightarrow\operatorname{Pic}(S)\longrightarrow\operatorname{Pic}(X)
\longrightarrow\operatorname{Pic}(X/S).
\]
When $X$ admits a section over $S$, the last homomorphism is surjective; that
is, there is then an isomorphism
\[
\operatorname{Pic}(X/S)\xrightarrow{\sim}
\operatorname{Pic}(X)/\operatorname{Pic}(S).
\]

\result{REMARK}{2.5} Let $f:X\longrightarrow S$ be a morphism such that
$\mathcal O_S\xrightarrow{\sim}f_*(\mathcal O_X)$, and let $g$ be a section of
$X$ over $S$. For an invertible Module $\mathcal L$ on $X$, call an isomorphism
$\mathcal O_S\xrightarrow{\sim}g^*(\mathcal L)$ a $g$-rigidification of
$\mathcal L$, and call an invertible Module $\mathcal L$ on $X$ equipped with a
$g$-rigidification a \emph{$g$-rigidified invertible Module}. Every
automorphism of such a structure is trivial, and $\operatorname{Pic}'(X/S)$ is
identified with the set of isomorphism classes of $g$-rigidified invertible
Modules on $X$\footnote{The printed text says “on $S$”; rigidification by a
section $g:S\to X$ here concerns the Module $\mathcal L$ on $X$.}. (This is the
fact that allowed descent theory to be used in proving 2.2(ii).) This gives a
new interpretation of $\operatorname{Pic}(X/S)$, at least when $f$ is also
quasi-compact and separated, so that
$\operatorname{Pic}(X/S)\simeq\operatorname{Pic}'(X/S)$ by 2.3.

\src{148}
\result{REMARK}{2.6} Let $f:X\longrightarrow S$ be as in 2.2, and let
$S'\longrightarrow S$ be faithfully flat and quasi-compact such that there is
an $S$-morphism $S'\longrightarrow X$, or equivalently such that
$X'=X\times_SS'$ has a section over $S'$. Put $S''=S'\times_SS'$ and
$X''=X\times_SS''$, and consider the exact sequence
\[
\operatorname{Pic}(X/S)\longrightarrow\operatorname{Pic}(X'/S')
\rightrightarrows\operatorname{Pic}(X''/S'').
\]
Applying 2.4 gives an exact sequence
\[
\operatorname{Pic}(X/S)\longrightarrow
\operatorname{Pic}(X')/\operatorname{Pic}(S')
\rightrightarrows
\operatorname{Pic}(X'')/\operatorname{Pic}(S''),
\]
and in particular every element of the relative Picard group already “comes
from” an element of $\operatorname{Pic}(X')$. This substantially simplifies
the description of the relative Picard group given in the preceding section,
and likewise that of the Picard functor of $X$ over $S$, since for every $T$
over $S$ the preceding argument applies to $X\times_ST/T$ and the morphism
$T'=S'\times_ST\longrightarrow T$. If, for example, $f$ itself is faithfully
flat, one may take $S'=X$. When $f$ is moreover of finite type (respectively
smooth, etc.), this permits one, in the description of the relative Picard
functor $T\longmapsto\operatorname{Pic}(X\times_ST/T)$, to restrict to base
changes $T'\longrightarrow T$ that are of finite type (respectively smooth,
etc.). When $f$ is projective and flat and $S$ is locally noetherian, one proves
that $S'\longrightarrow S$ above may be chosen so that $S'$ is the disjoint
union of flat coverings $S'_i$ of open subsets $S_i$ covering $S$; if $f$ is
also separable, $S'_i$ may be taken étale over $S_i$.

\fsection{3}{The main existence theorem: statement}

There is not, even conjecturally, an existence statement for Picard preschemes
that encompasses all known cases. A “practically necessary” condition, if one
may say so, is that $f:X\longrightarrow S$ be proper, ensuring essential
finiteness properties, and flat. These conditions are not sufficient, even if
$S$ is the spectrum of the algebra of dual numbers $k[t]/(t^2)$ over a field
$k$ (say the field $\mathbf C$ of complex numbers) and $X$ has dimension 1. At
the time of writing, the most important existence theorems for the Picard
prescheme are deduced from the following theorem.

\result{THEOREM}{3.1} Let $f:X\longrightarrow S$ be a morphism of locally
noetherian preschemes. Suppose:

\noindent (i) $f$ is projective;

\src{149}
\noindent (ii) $f$ is flat;

\noindent (iii) the geometric fibers of $f$ are integral.

Under these conditions, $\operatorname{Pic}_{X/S}$ exists.

The proof, which will be sketched in the next two sections, will show at the
same time the following. Let $\xi$ be the section of
$\operatorname{Pic}_{X/S}$ corresponding to a very ample sheaf
$\mathcal O_X(1)$ on $X/S$, that is, one induced by a projective immersion
$X\longrightarrow\mathbf P(\mathcal E)$. There is then an open subset $U$ of
$\operatorname{Pic}_{X/S}$, a disjoint union of open subsets quasi-projective
over $S$, such that $U$ is stable under translation by $\xi$, and
$\operatorname{Pic}_{X/S}$ is the increasing union of the open subsets
$U-n\xi$, all isomorphic to $U$. In particular, under the conditions of 3.1,
$\operatorname{Pic}_{X/S}$ is separated over $S$.

\result{REMARK}{3.2} Examples show (for instance with $S$ the spectrum of a
discrete valuation ring and $X$ of relative dimension 1 over $S$) that if
hypothesis (iii) in 3.1 is omitted and replaced by the weaker hypothesis that,
for every $s\in S$, the homomorphism
\[
k(s)\longrightarrow H^0(X_s,\mathcal O_{X_s})
\]
is an isomorphism, then $\operatorname{Pic}_{X/S}$ need not be separated over
$S$. This happens both when the geometric fibers of $f$ are reduced but an
integral geometric generic fiber “bursts” under specialization into two
irreducible components, and when the geometric fibers of $f$ are irreducible
but an integral geometric generic fiber specializes to a “multiple fiber.”
The first case occurs, for example, with a conic degenerating into two
intersecting lines; D. Mumford supplied an example of the second, with an
elliptic curve degenerating into a double elliptic curve. These examples work
in every characteristic.

\result{REMARK}{3.3} Under the conditions of 3.1, I do not know whether
$\operatorname{Pic}_{X/S}$ is a disjoint union of open subsets of finite type,
and hence quasi-projective, over $S$.\footnote{The 1962 erratum reports that
D. Mumford answered this question affirmatively by means of the theory of
Picard schemes developed here.} The Hilbert polynomials
$Q\in\mathbf Q[t]$, as in the case of Hilbert schemes ([4], IV), give a
decomposition of $\operatorname{Pic}_{X/S}$ into a disjoint union of open
subsets $\operatorname{Pic}^{Q}_{X/S}$, and it seems plausible that the latter
are quasi-projective over $S$; this will at least be seen in the next exposé
when $f$ is smooth. Notice that if hypothesis (i) is replaced by the hypothesis
that $X$ is locally projective over $S$ (which suffices for 3.1, since the
existence of $\operatorname{Pic}_{X/S}$ is visibly local on $S$), it is, on the
other hand, easy

\src{150}
to give examples in which $\operatorname{Pic}_{X/S}$ contains connected
components that are not of finite type over $S$. For example, let $X_0$ be a
nonsingular projective algebraic variety over an algebraically closed field
$k$, equipped with an automorphism $u$ and an element $\xi$ of the
Néron--Severi group of $X_0$ such that the $u^n(\xi)$ are pairwise distinct.
One may take $X_0$ to be the product of an elliptic curve $E$ with itself and
$u$ to be the automorphism
\[
(x,y)\longmapsto(x,y+x)
\]
of $E\times E$. Let $S$ be the union of two nonsingular irreducible curves
meeting at two points $a$ and $b$. There is a connected principal
$\mathbf Z$-covering $P$ over $S$; using the actions of $\mathbf Z$ on $X_0$
defined by $u$, one obtains an associated fiber space over $S$ with fiber
$X_0$, trivial over $S-a$ and $S-b$ and, in the particular case considered,
in fact an abelian scheme over $S$. It is readily seen that
$\operatorname{Pic}_{X/S}$, which is also the fiber space associated with $P$
and the actions of $\mathbf Z$ on $\operatorname{Pic}_{X_0/k}$ through $u$,
contains a connected component isomorphic to
\[
P\times\operatorname{Pic}^{0}_{X_0/k},
\]
where $\operatorname{Pic}^{0}$ denotes the connected component of the identity
in $\operatorname{Pic}$; this component is not of finite type over $S$.
[Analogous phenomena also occur in various cases of Picard preschemes
nonseparated over $S$, as considered in 3.2.]

\fsection{4}{Relative Cartier divisors and projective bundles}

We shall need only positive divisors and will omit that additional qualifier
throughout this section.

Let $X$ be a prescheme. A \emph{Cartier divisor} on $X$, or simply a divisor on
$X$, is a closed subprescheme $D$ of $X$ defined by an Ideal $\mathcal J$ that
is an invertible Module, that is, locally generated by a section that is not a
zero divisor in $\mathcal O_X$. With $D$ we associate the invertible Module
\[
\mathcal L(D)=\mathcal J^{-1}.
\]
The canonical injection $\mathcal J\longrightarrow\mathcal O_X$ gives a
canonical homomorphism
\[
s_D:\mathcal O_X\longrightarrow\mathcal J^{-1}=\mathcal L(D),
\qquad\text{that is,}\quad s_D\in\Gamma(X,\mathcal L(D)).
\]
Moreover, giving a divisor is essentially equivalent to giving an invertible
Module $\mathcal L$ on $X$ together with a section $s$ that is nowhere a zero
divisor, by associating conversely with such a pair $(\mathcal L,s)$ the
“divisor” of $s$, denoted $\operatorname{div}(s)$. For a fixed invertible
$\mathcal L$ on $X$, the set of

\src{151}
divisors $D$ defining $\mathcal L$ is in bijection with the quotient set
\[
\Gamma(X,\mathcal L)^*/\Gamma(X,\mathcal O_X^*),
\]
where $\Gamma(X,\mathcal L)^*$ is the subset of $\Gamma(X,\mathcal L)$ formed
by sections that are nowhere zero divisors.

Now suppose given a locally finite-type morphism $f:X\longrightarrow S$, with
$S$ locally noetherian for simplicity. Let $\mathcal J$ be a coherent Ideal on
$X$, let $D$ be the subscheme of $X$ that it defines, let $x\in X$, and put
$s=f(x)$. The following conditions are equivalent:

\noindent (i) $\mathcal J$ is invertible at $x$ (that is, $\mathcal J_x$ is
generated by a regular element of $\mathcal O_{X,x}$), and $D$ is flat over
$S$ at $x$;

\noindent (ii) $X$ and $D$ are flat over $S$ at $x$, and $D_s$ is a Cartier
divisor on the fiber $X_s$ at $x$;

\noindent (iii) $X$ is flat over $S$ at $x$, and $\mathcal J_x$ is generated
by an element $f_x$ inducing on $X_s$ a germ that is not a zero divisor.

One then says that $D$ is a \emph{relative Cartier divisor} on $X/S$ (or simply
a relative divisor on $X/S$) at the point in question. Condition (i) shows
that $D$ is also a relative divisor at neighboring points of $x$. Thus, if $X$
and $D$ are flat over $S$ and $D$ is proper over $S$, the set of $s\in S$ such
that $D_s$ is a Cartier divisor in $X_s$ (that is, such that $D$ is a relative
Cartier divisor at the points of $X_s$) is open in $S$. Furthermore, the
preceding definition was made precisely so that relative Cartier divisors are
stable under arbitrary base change $S'\longrightarrow S$. Let
$\operatorname{Div}(X/S)$ be the set of relative divisors on $X/S$, and define
a contravariant functor in $S'$ over $S$ by
\[
\operatorname{Div}_{X/S}(S')=
\operatorname{Div}(X\times_SS'/S').
\]

Suppose $X$ flat and proper over $S$. By characterization (ii) of relative
Cartier divisors, $\operatorname{Div}_{X/S}$ may be regarded as a subfunctor of
the functor $\operatorname{Hilb}_{X/S}$ defined in [4], IV, and the inclusion
\[
\operatorname{Div}_{X/S}\longrightarrow\operatorname{Hilb}_{X/S}
\]
is “representable by open immersions” (cf. [5], IV, 3.13), by the preceding
remarks. The main existence theorem of [4], IV, therefore gives:

\src{152}
\result{PROPOSITION}{4.1} Suppose $f:X\longrightarrow S$ is projective and
flat. Then the functor $\operatorname{Div}_{X/S}$ is representable; more
precisely, it is represented by an open subset of
$\operatorname{Hilb}_{X/S}$.

Using, for a fixed very ample sheaf $\mathcal O_X(1)$ on $X/S$, the canonical
decomposition of $\operatorname{Hilb}_{X/S}$ into the open subsets
$\operatorname{Hilb}^{Q}_{X/S}$ corresponding to the Hilbert polynomials
$Q\in\mathbf Q[t]$, one also obtains the analogous decomposition
\[
\operatorname{Div}_{X/S}=
\coprod_{Q\in\mathbf Q[t]}\operatorname{Div}^{Q}_{X/S}
\]
as a disjoint union of open subsets quasi-projective over $S$.

The assignment $D\longmapsto\mathcal L(D)$ also gives a functorial
homomorphism
\begin{equation}
\operatorname{Div}_{X/S}\longrightarrow\operatorname{Pic}_{X/S}
\tag{+}
\end{equation}
which we now study; under fairly general conditions it will prove relatively
representable ([5], IV, 3). Begin with an element $\xi$ of
$\operatorname{Pic}_{X/S}(S')$, take $S'=S$ to simplify notation, and show that
the corresponding subfunctor of $\operatorname{Div}_{X/S}$ is representable.
First consider the case in which $\xi$ is defined by an invertible Module
$\mathcal L$ on $X$. Suppose that $X$ is proper and flat over $S$ and that its
geometric fibers over $S$ are integral. This also implies ([2], III, par. 7)
that $\mathcal O_S\xrightarrow{\sim}f_*(\mathcal O_X)$ and that this relation
remains valid after every base change $S'\longrightarrow S$. The relative
Cartier divisors $D$ on $X/S$ for which $\mathcal L(D)$ and $\mathcal L$ define
the same element of
$\operatorname{Pic}(X/S)=\operatorname{Pic}_{X/S}(S)$, equivalently by 2.4 for
which $\mathcal L(D)$ and $\mathcal L$ are isomorphic locally over $S$, are in
bijection with the sections of the quotient sheaf
$f_*(\mathcal L)^*/\mathcal O_S^*$. This correspondence is compatible with
base change. General considerations of the “Künneth” type from the cited
passage (see also [6]) show that the properness of $X/S$ and the flatness of
$\mathcal L$ over $S$ imply the existence of a coherent Module $\mathcal Q$ on
$S$, defined up to unique isomorphism, and a sheaf isomorphism
\[
f_*(\mathcal L)\xrightarrow{\sim}
\operatorname{Hom}_{\mathcal O_S}(\mathcal Q,\mathcal O_S).
\]
\footnote{The printed text places $\mathcal O_X$ under
$\operatorname{Hom}$; since $\mathcal Q$ is a Module on $S$, its dual is taken
over $\mathcal O_S$.} The formation of $\mathcal Q$ is compatible with base
change. Here $f_*(\mathcal L)^*$ denotes the subsheaf of sets of
$f_*(\mathcal L)$ whose sections over $U$ are the sections of $\mathcal L$ over
$f^{-1}(U)$ that define relative Cartier divisors on

\src{153}
$f^{-1}(U)/U$, that is, that induce sections which are not zero divisors on the
$X_s$ for $s\in U$. Since the fibers $X_s$ are integral, this simply means that
the induced sections on the fibers $X_s$ are not identically zero or,
equivalently, in terms of local homomorphisms
$\mathcal Q\longrightarrow\mathcal O_S$, that these homomorphisms are
surjective, by Nakayama's lemma. Thus the set of sections of
$f_*(\mathcal L)^*/\mathcal O_S^*$ is in bijection with the set of invertible
quotient Modules of $\mathcal Q$, or, by the definition of the projective
bundle $\mathbf P(\mathcal Q)$ associated with the coherent Module
$\mathcal Q$ (cf. [5], V, 2), with the set of sections of
$\mathbf P(\mathcal Q)$ over $S$. This description is compatible with inverse
images, and hence yields the following theorem.

\result{THEOREM}{4.3} Let $f:X\longrightarrow S$ be a proper flat morphism with
integral geometric fibers, where $S$ is locally noetherian, and let
$\mathcal L$ be an invertible Module on $X$. For every $S'$ over $S$, let
$T(S')$ be the set of relative divisors $D$ on $X\times_SS'/S'$ such that
$\mathcal L(D)$ is isomorphic to
$\mathcal L\otimes_{\mathcal O_S}\mathcal O_{S'}$ locally over $S'$, that is,
such that $\mathcal L(D)$ and
$\mathcal L\otimes_{\mathcal O_S}\mathcal O_{S'}$ define the same element of
$\operatorname{Pic}(X\times_SS'/S')$. Then there is a coherent Module
$\mathcal Q$ on $S$, determined up to unique isomorphism, such that the functor
$T$ is represented by the projective bundle $\mathbf P(\mathcal Q)$.

\result{COROLLARY}{4.4} If $f$ is moreover projective, then the functorial
homomorphism
$\operatorname{Div}_{X/S}\longrightarrow\operatorname{Pic}_{X/S}$ is
representable by projective morphisms.

Indeed, when $X$ admits a section (respectively, locally a section) over $S$,
the preceding homomorphism is represented by projective bundles (respectively,
local projective bundles), by 4.3 and 2.1. When $f$ is quasi-projective, one
readily reduces to the preceding case by a descent argument using finite flat
quasi-sections of $X$ locally over $S$.

\result{REMARK}{4.5} Under the conditions of 4.3, the Module $\mathcal Q$ on
$S$ is not generally locally free, as expressed by the fact that the dimension
of the reduced fibers of $\mathcal Q$, that is, the dimension of
$H^0(X_s,\mathcal L_s)$ for varying $s\in S$, may jump. Given a coherent
Module $\mathcal Q$ on a locally noetherian prescheme $S$, one readily checks
that for a fixed $s\in S$, $\mathcal Q$ is free at $s$ if and only if
$\mathbf P(\mathcal Q)$ is flat over $S$ at the points above $s$; in that case
it is even smooth over $S$ at those points. Here, with $\mathcal Q$ defined in
terms of $\mathcal L$ as above, this also means that formation of the direct
image $f_*(\mathcal L)$ “commutes with base change

\src{154}
near $s$,” or equivalently
that $f_*(\mathcal L)_s\longrightarrow H^0(X_s,\mathcal L_s)$ is surjective.
This holds, for example, if $H^1(X_s,\mathcal L_s)=0$. Subject to the existence
of the preschemes involved, these criteria apply in particular to the
universal situation
$\operatorname{Div}_{X/S}\longrightarrow\operatorname{Pic}_{X/S}$ and give a
necessary and sufficient condition, respectively a sufficient condition, for
this morphism to be smooth at a given point of $\operatorname{Div}_{X/S}$.

\fsection{5}{Proof of the main existence theorem}

Under the conditions of 3.1, choose a very ample Module $\mathcal O_X(1)$ on
$X/S$, and let $\xi$ be the corresponding element of
$\operatorname{Pic}(X/S)$. For brevity put
\[
P(S')=\operatorname{Pic}(X\times_SS'/S'),
\]
and suppose first, for simplicity, that $X/S$ admits a section. Let $P^+(S')$
be the subset of $P(S')$ consisting of the classes of invertible $\mathcal L$
on $X\times_SS'$ such that
\[
R^if'_*(\mathcal L(n))=0\quad\text{for $i>0$ and every $n\geq0$},
\qquad
f'_*(\mathcal L(n))\neq0\quad\text{for every $n\geq0$}.
\]
These conditions are stable under base change and hence define a subfunctor
$P^+$ of $P$, plainly stable under translation by $\xi$. Using Serre's
“Theorems A and B” ([2], III, 2) and the generalities in [5], IV, 5, one readily
sees that $P$ is representable if and only if $P^+$ is; in that case $P^+$ is
represented by an open subset $U$ of the prescheme
$\operatorname{Pic}_{X/S}$ representing $P$, and the latter is the increasing
union of the open subsets $U-n\xi$.

For brevity put $D=\operatorname{Div}_{X/S}$, and let $D^+$ be the inverse
image of $P^+$ under the canonical morphism $D\longrightarrow P$. Thus there
is a morphism
\begin{equation}
D^+\longrightarrow P^+,
\tag{+}
\end{equation}
and $D^+$ is already known to be represented by an open subset $D^+$ of the
prescheme $D=\operatorname{Div}_{X/S}$, whose existence is guaranteed by 4.1.
It follows readily from 4.3 that (+) is represented by smooth projective
surjective morphisms, more precisely by projective bundles associated with
everywhere nonzero locally free Modules. Indeed, when $\mathcal L$ on
$X\times_SS'$ is as at the beginning of this section, $f'_*(\mathcal L)$ is a
nonzero locally free Module whose formation commutes with base change; in the
notation of 4.3, $\mathcal Q$ is then isomorphic to the dual of
$f'_*(\mathcal L)$. We shall deduce the representability of $P^+$ from the
general criterion [5], IV, 4.7. In the cited passage, take $\mathfrak C$ to be
the class of faithfully

\src{155}
flat quasi-compact prescheme morphisms, which are effective epimorphisms by
[4], I B. Condition (a) of the cited passage, namely that (+) is representable
by morphisms belonging to $\mathfrak C$, has just been verified. So has
condition (b), which says that the functor $P^+$ is compatible with faithfully
flat quasi-compact descent; this is immediate. It remains to prove condition
(c) of that passage: the equivalence relation $R$ in the prescheme $D^+$
deduced from the $\mathfrak C$-representable morphism (+) is
$\mathfrak C$-effective, that is, it is effective and
$D^+\longrightarrow D^+/R$ belongs to $\mathfrak C$. First observe that the
open subsets $D^{+Q}$ of $D^+$ corresponding to the various virtual Hilbert
polynomials $Q\in\mathbf Q[t]$ are stable under $R$, since the fibers of $R$
are connected. It therefore suffices to show, for every $Q$, that the
equivalence relation $R^Q$ induced in $D^{+Q}$ is $\mathfrak C$-effective.
Now $D^{+Q}$ is quasi-projective, while the equivalence relation $R^Q$ is
projective and flat. Thus the hypotheses of [4], III, 6.1 apply and give the
required result.

In the general case in which $X/S$ does not necessarily admit a section, one
readily reduces to the preceding case by descent, or repeats the preceding
proof with the required modification in the definition of $P^+$.

\result{REMARKS}{5.1} The method used is essentially Matsusaka's method for the
projective construction of Picard varieties. The result from [4], III invoked
for passage to the effective quotient can also be deduced readily from the
existence theorem for Hilbert schemes (see, for example, [6]); classically,
these quotients were constructed with Chow coordinates. Notice that formation
of the open subset $\operatorname{Pic}^{+}_{X/S}$ of
$\operatorname{Pic}_{X/S}$ and its decomposition into open subsets
$\operatorname{Pic}^{+Q}_{X/S}$ quasi-projective over $S$, according to the
Hilbert polynomials of the divisors defining the invertible Modules in
question, is compatible with base change, which permits the descent argument.

\result{REMARKS}{5.2} It is possible that $\operatorname{Pic}_{X/S}$ exists
whenever $f:X\longrightarrow S$ is proper and flat and the homomorphisms
\[
k(s)\longrightarrow H^0(X_s,\mathcal O_{X_s})\qquad(s\in S)
\]
are isomorphisms; this last condition then also means that
$\mathcal O_S\xrightarrow{\sim}f_*(\mathcal O_X)$ and that this relation
remains valid after every base change $S'\longrightarrow S$. This is at least
what is proved for analytic spaces when $f$ is also projective, by a different
method, Chow's method unless I am mistaken, explained in [5], IX,
3.1.\footnote{The 1962 erratum states that the existence conjecture advanced
at the beginning of this remark is false, although Mumford's methods prove a
slightly weaker theorem.} In that method, passage to the

\src{156}
quotient by a proper flat equivalence relation in an open subset of the scheme
of divisors is replaced by passage to the quotient by the projective group in
the scheme of immersions of $X$ in $\mathbf P^r_S$. This method can probably
be adapted to schemes using Mumford's results on passage to the quotient by
the projective group [8]. For the present, a written proof exists only when
$X$ has “many local sections” over $S$, for example when $X$ is separable over
a complete local ring with algebraically closed residue
field.\footnote{The printed text omits “with algebraically closed residue
field”; the erratum says that Mumford's counterexample makes this restriction
indispensable.} In principle, this method would be more general, since
it also gives the existence of Picard preschemes in cases in which they are
nonseparated, where the first method must therefore fail. (Technically, the
difficulty is that when the geometric fibers of $f$ are no longer integral,
the functor considered in 4.3 is represented not by the projective bundle
$\mathbf P(\mathcal Q)$ itself but by an open subset of it, leading to the
delicate problem of taking the quotient by a flat but nonproper equivalence
relation.)

\result{REMARK}{5.3} Notice that the proof given here uses neither a prior
construction of the Jacobians of curves or families of curves, nor the theory
of abelian varieties or abelian schemes. It therefore differs essentially
from traditional treatments such as Lang's book [7] or Chevalley's account
[1], which follow the path sketched by A. Weil. Even for the Jacobians of
nonsingular curves over an algebraically closed field, say the complex
numbers, the construction given here is the only known one providing the very
strong properties taken as the definition in Section 1: essentially
Chevalley's, but allowing “parameter varieties” with nilpotent elements. That
the construction of Picard schemes should precede rather than follow the
theory of abelian varieties is clear a priori, since in general Picard schemes
are neither abelian schemes nor reducible to them. This is already seen for
singular curves over an algebraically closed field, where Rosenlicht's
“generalized Jacobians,” which are not abelian varieties, occur. Moreover, the
theory of abelian varieties, and more generally of abelian schemes, becomes
much simpler once a general theory of Picard schemes is available. In
particular, duality theory for abelian schemes, notably results of Cartier
type, then becomes nearly formal (see, for example, [10]).

\src{157}
\result{REMARKS}{5.4} Igusa's “compatibility principle” for the Jacobian of a
curve degenerating to a singular curve can be properly understood only as an
existence theorem for the Picard scheme of a relative curve scheme $X/S$, not
necessarily smooth over $S$. It is therefore a special case of the main
existence theorem 3.1 when the specialized curve is integral, that is, in
classical language, irreducible of multiplicity 1. At present, the case of a
reducible special curve, even when its components have multiplicity 1, that
is, when the special curve is separable over the residue field, lies outside
the known existence theorems, except over a complete discrete valuation ring
with algebraically closed residue field; compare 5.2. This existence question
will certainly arise in a scheme-theoretic geometric construction of the
Baily--Satake “compactifications” of the moduli schemes of curves of genus
$g$. (This compactification is currently known only for $g=1$, through
Igusa's work.)

\fsection{6}{Relative existence theorems}

We sketch here some useful cases in which the existence of certain Picard
schemes implies that of others, allowing various further existence theorems
to be deduced from the main theorem 3.1.

\result{PROPOSITION}{6.1} Let $f:X\longrightarrow S$ be a flat projective
morphism such that in the Stein factorization $f=f''f'$, the morphism
$f':X\longrightarrow S'$ is flat with integral geometric fibers, and hence
satisfies the hypotheses of 3.1, while the finite morphism
$f'':S'\longrightarrow S$ is flat. Then $\operatorname{Pic}_{X/S}$ exists and,
with the notation introduced in [4], II, C 2, there is a canonical isomorphism
\[
\operatorname{Pic}_{X/S}\xrightarrow{\sim}
\Pi_{S'/S}\operatorname{Pic}_{X/S'}.
\]

To prove this, one first establishes an isomorphism of functors
\[
\operatorname{Pic}_{X/S}\xrightarrow{\sim}
\Pi_{S'/S}\operatorname{Pic}_{X/S'},
\]
uses 3.1, which implies that $\operatorname{Pic}_{X/S'}$ is representable, and
then uses the structure of $\operatorname{Pic}_{X/S'}$ described in Section 3,
which implies that every finite subset of a fiber of
$\operatorname{Pic}_{X/S'}$ over $S$ is contained in an affine open subset, to
obtain the existence of $\Pi_{S'/S}\operatorname{Pic}_{X/S'}$.

\src{158}
For example, when $X$ is a scheme that is the disjoint union of schemes $X_i$
over $S$ satisfying the conditions of 3.1, statement 6.1 reduces to the trivial
statement
\[
\operatorname{Pic}_{X/S}\xrightarrow{\sim}
\prod_i\operatorname{Pic}_{X_i/S}.
\]

\result{COROLLARY}{6.2} Let $f:X\longrightarrow S$ be a projective flat
morphism with locally integral geometric fibers (for example, a projective
normal morphism). Then $\operatorname{Pic}_{X/S}$ exists.

In this case, moreover, $S'$ is an étale covering of $S$; this holds whenever
$f$ is separable, that is, flat with geometrically reduced fibers. One checks
that the structure theorem stated in Section 3 for
$\operatorname{Pic}_{X/S}$ remains valid, by virtue of the analogous structure
of $\operatorname{Pic}_{X/S'}$.

On the other hand, a descent procedure gives a relative existence theorem
whose reach depends on the solution of the nonflat descent questions raised in
[4], I, A (c). We spell out only the following particular case:

\result{PROPOSITION}{6.3} Let $f:X\longrightarrow S$ be a proper morphism,
and let $X_1$ and $X_2$ be two subschemes of $X$, flat over $S$, defined by two
coherent Ideals $\mathcal I_1$ and $\mathcal I_2$ such that
$\mathcal I_1\cap\mathcal I_2=(0)$ and
$\mathcal O_X/(\mathcal I_1+\mathcal I_2)$ is flat over $S$; that is, the
subscheme of $X$ that is the supremum of $X_1$ and $X_2$ is $X$, while their
infimum $Z$ is flat over $S$.\footnote{The final letter is printed “$X$,” but
the immediately preceding condition concerns the flatness of
$\mathcal O_X/(\mathcal I_1+\mathcal I_2)$ over $S$ and therefore defines $Z$
to be flat over $S$.} Suppose, in addition, that for every $s\in S$ the
homomorphisms
\[
k(s)\longrightarrow H^0(X_{i,s},\mathcal O_{X_{i,s}}),\qquad i=1,2,
\]
are bijective. Then the natural homomorphism of functors
\[
\operatorname{Pic}_{X/S}\longrightarrow
\operatorname{Pic}_{X_1/S}\times\operatorname{Pic}_{X_2/S}
\]
is representable by affine morphisms. Consequently, if
$\operatorname{Pic}_{X_1/S}$ and $\operatorname{Pic}_{X_2/S}$ exist, then so
does $\operatorname{Pic}_{X/S}$, and the canonical morphism
\[
\operatorname{Pic}_{X/S}\longrightarrow
\operatorname{Pic}_{X_1/S}\times\operatorname{Pic}_{X_2/S}
\]
is affine.

By faithfully flat descent one reduces to the case in which $Z$ admits a
section over $S$, thereby defining sections of $X$, $X_1$, and $X_2$ over $S$
and allowing the automorphisms in the structures under consideration to be
eliminated as explained in 2.5. The proof then consists in observing that the
data of an invertible

\src{159}
“rigidified” Module $\mathcal L$ on $X$ are equivalent to the data of a triple
$(\mathcal L_1,\mathcal L_2,u)$, where $\mathcal L_i$ is a “rigidified”
invertible Module on $X_i$ and $u$ is an isomorphism from
$\mathcal L_1|Z$ to $\mathcal L_2|Z$ compatible with the rigidifications. One
need only verify that, for fixed $\mathcal L_1$ and $\mathcal L_2$, the datum
of $u$ is expressed by a section of a suitable scheme over $S$ that is affine
over $S$; this is easy. One readily deduces from 6.3:

\result{COROLLARY}{6.4} Let $X$ be a proper separable scheme over a field $k$,
and let $X_i$ be the irreducible components of $X$. If the
$\operatorname{Pic}_{X_i/k}$ exist, then so does
$\operatorname{Pic}_{X/k}$, and the canonical morphism
\[
\operatorname{Pic}_{X/k}\longrightarrow\prod_i\operatorname{Pic}_{X_i/k}
\]
is affine.

Combined with 6.2, this proves, for example, the existence of
$\operatorname{Pic}_{X/k}$ whenever $X$ is a projective separable scheme over
a field $k$. When $X$ is no longer separable over $k$, there is also a
reduction result, using an argument due to Oort [11]. The method also applies
over an arbitrary base scheme, a case useful, for example, in proving in the
next exposé the finiteness result announced in 3.3. To avoid an overly
technical statement, let us restrict ourselves to a base field:

\result{PROPOSITION}{6.5} Let $X$ be a proper scheme over a field $k$, and let
$X_0$ be a subscheme with the same underlying set, hence defined by a
nilpotent Ideal on $X$. Then the functorial morphism
\[
\operatorname{Pic}_{X/k}\longrightarrow\operatorname{Pic}_{X_0/k}
\]
is representable by affine morphisms. In particular, if
$\operatorname{Pic}_{X_0/k}$ exists, then so does
$\operatorname{Pic}_{X/k}$, and the morphism
$\operatorname{Pic}_{X/k}\longrightarrow\operatorname{Pic}_{X_0/k}$ is
affine.

Combining this with 6.4, one readily concludes:

\result{COROLLARY}{6.6} Let $X$ be a projective prescheme over a field $k$.
Then $\operatorname{Pic}_{X/k}$ exists.

\result{REMARK}{6.6} It is extremely plausible that
$\operatorname{Pic}_{X/k}$ exists for every proper scheme $X$ over a field
$k$.\footnote{The 1962 erratum reports that Murre had just answered this
question affirmatively.} The preceding results reduce this question to the
case in which $k$ is algebraically closed and $X$ is integral. There is then
an integral scheme $X'$, projective over $k$, and a dominant morphism
$g:X'\longrightarrow X$ by Chow's lemma. It would therefore suffice to show
that the

\src{160}
corresponding functorial morphism
\[
\operatorname{Pic}_{X/k}\longrightarrow\operatorname{Pic}_{X'/k}
\]
is representable, and no doubt representable by affine morphisms, since
$\operatorname{Pic}_{X'/k}$ is already known to be representable. This raises
nonflat descent questions that remain unresolved for the present. Notice that
if one restricts the functor $\operatorname{Pic}_{X/k}$ to reduced preschemes,
with $X$ proper and integral over an algebraically closed field $k$, one does
obtain a representable functor, as Chevalley [1] proved when $X$ is normal and
Seshadri [12] proved in general by a descent method. In our notation, however,
the scheme constructed by these authors is not $\operatorname{Pic}_{X/k}$ but
$(\operatorname{Pic}_{X/k})_{\mathrm{red}}$, the reduced scheme corresponding
to $\operatorname{Pic}_{X/k}$.

\result{REMARK}{6.7} More generally, let $f:X'\longrightarrow X$ be a
surjective morphism of proper preschemes over $k$. Nonflat descent
considerations lead one to conjecture that
\[
\operatorname{Pic}_{X/k}\longrightarrow\operatorname{Pic}_{X'/k}
\]
is an affine morphism.\footnote{The 1962 erratum reports that this question has
an affirmative answer and refers to the last paragraph of the Comments on the
next exposé.} This would imply in particular, after dividing by the connected
components of the identity elements, that the corresponding homomorphism on
Néron--Severi groups is injective modulo torsion. This can be checked by
intersection theory when $X$ and $X'$ are nonsingular. The answer does not
seem to be known in any other case.

\result{REMARK}{6.8} Contrary to what one might think, Picard schemes of
algebraic schemes with nilpotent elements are useful, and even indispensable,
in various questions. Thus, when $X$ is a projective scheme, say smooth over
$k$,\footnote{The printed text reads “regular”; the erratum prescribes “smooth
over $k$.”} and $Y$ is a hyperplane section, one must consider the
“infinitesimal neighborhoods” $X_n$ of $Y$ of all orders and the Picard
schemes $\operatorname{Pic}_{X_n/k}$. When $X$ is irreducible of dimension
$\geq4$ (respectively $\geq3$), the canonical morphism
\[
\operatorname{Pic}_{X/k}\longrightarrow\operatorname{Pic}_{X_n/k}
\qquad(n\text{ large})
\]
is an isomorphism (respectively induces an isomorphism on the inverse images
of the torsion subgroups of Néron--Severi), a result that will be useful in the
qualitative study of Picard schemes in the next exposé. Likewise, considering
the Picard schemes of certain curves with nilpotent elements, together with
the fundamental theorems of formal geometry [3], makes it possible to verify,
in the equal-characteristic case, a conjecture of Mumford: for every complete
normal noetherian local ring $A$ of dimension 2, the divisor class group of
$A$ may be regarded as the set of $k$-rational points of an

\src{161}
algebraic group $G$ over the residue field $k$; moreover, $G$ is canonically
determined once a coefficient field in $A$ is chosen. When $A$ has arbitrary
dimension, it is plausible that there exists a pro-algebraic group over $k$
playing the same role as $G$ above. In the case in which
$\operatorname{Spec}(A)$ can be “desingularized,” it is constructed as an
inverse limit of Picard schemes of suitable projective schemes over $k$ with
nilpotent elements.

\section*{BIBLIOGRAPHY}
\addcontentsline{toc}{section}{Bibliography}

\begin{description}
\item[{[1]}] CHEVALLEY (Claude). -- Sur la théorie de Picard,
\emph{Amer. J. of Math.}, vol. 82, 1960, pp. 435--490.

\item[{[2]}] GROTHENDIECK (A.) and DIEUDONNÉ (J.). -- Éléments de géométrie
algébrique, Chapitre I et suivants. -- Paris, Presses universitaires de France,
1960--1961 (Institut des Hautes Études Scientifiques, Publications
mathématiques 4, 8, \ldots).

\item[{[3]}] GROTHENDIECK (Alexander). -- Géométrie formelle et géométrie
algébrique, \emph{Séminaire Bourbaki}, vol. 11, 1958/59, no. 182, 28 pp.

\item[{[4]}] GROTHENDIECK (Alexander). -- Techniques de descente et théorèmes
d'existence en géométrie algébrique, I--IV, \emph{Séminaire Bourbaki}, vol. 12,
1959/60, no. 190, 29 pp. and no. 195, 22 pp.; vol. 13, 1960/61, no. 212,
20 pp. and no. 221, 28 pp.

\item[{[5]}] GROTHENDIECK (Alexander). -- Techniques de construction en
géométrie analytique, I--X, \emph{Séminaire Cartan}, vol. 13, 1960/61:
Déformation des structures analytiques complexes, nos. 7--17.

\item[{[6]}] GROTHENDIECK (Alexander). -- Séminaire de géométrie algébrique,
edited by Lichtenbaum, Harvard University, 1961 (forthcoming).

\item[{[7]}] LANG (Serge). -- \emph{Abelian varieties}. -- New York,
Interscience Publishers, 1959 (Interscience Tracts in pure and applied
Mathematics, 7).

\item[{[8]}] MUMFORD (D.). -- An elementary theorem in geometric invariant
theory, \emph{Bull. Amer. math. Soc.}, vol. 67, 1961, pp. 483--486.

\item[{[9]}] MUMFORD (D.). -- The topology of normal singularities of an
algebraic surface and a criterion for simplicity. -- Paris, Presses
universitaires de France, 1961 (Institut des Hautes Études Scientifiques,
Publications mathématiques, 9).

\item[{[10]}] MUMFORD (D.) and TATE (J.). -- Séminaire de géométrie
algébrique, Harvard University, Spring Term 1962 (forthcoming).

\item[{[11]}] OORT (F.). -- Sur le schéma de Picard,
\emph{Bull. Soc. math. France}, vol. 90, 1962 (forthcoming).

\item[{[12]}] SESHADRI (C. S.). -- Thesis, forthcoming in
\emph{Annali di Matematica}.
\end{description}

\par\smallskip
\begingroup\footnotesize
\endgroup
